\section{Integral class pairings and the initial period bound}
\label{cp-section}

The initial period valuation has an elementary proof once the
actual polynomial lattices and their class coordinates have been
identified. The period is a finite class sum of the normalized
bilinear Fock pairing. Its only denominator at this stage is the
order of a finite stabilizer. We give the argument with arbitrary
finite stabilizers, so it also applies to the row classes with
additional level characters. The stronger divisibility of the
iterated columns is a separate assertion.

\subsection{Equivariant functions and integral pairings}

Let $K$ be a number field, $\mathcal O$ its ring of integers, and
let a group $G$ act on a set $\mathcal D$ with finitely many orbits.
Choose representatives $d_c$, indexed by a finite set $C$, and
suppose their stabilizers $\Gamma_c$ are finite. Put
$n_c=|\Gamma_c|\ge1$. Let $V^-$ and $V^+$ be finite-dimensional
$K$-spaces with representations $\rho^-$ and $\rho^+$ of $G$.
Suppose the nondegenerate bilinear pairing
\[
                    B:V^-\times V^+\longrightarrow K
\]
satisfies
\begin{equation}\label{cp-invariance}
 B(\rho^-(g)x,\rho^+(g)y)=B(x,y)\qquad(g\in G).
\end{equation}
For each $d\in\mathcal D$ let $M_d^\pm$ be full
$\mathcal O$-lattices, with
\[
 M_{gd}^\pm=\rho^\pm(g)M_d^\pm,
 \qquad B(M_d^-,M_d^+)\subseteq\mathcal O.
\]
Define the spaces and their integral lattices by
\[
 \begin{split}
 \mathcal A^\pm_K
   &=\{f:\mathcal D\to V^\pm:
                         f(gd)=\rho^\pm(g)f(d)\},\\
 \mathcal A^\pm_{\mathcal O}
   &=\{f\in\mathcal A^\pm_K:f(d)\in M_d^\pm
                                      \text{ for every }d\}.
 \end{split}
\]
Evaluation at the $d_c$ identifies these objects with
\[
 \bigoplus_c(V^\pm)^{\Gamma_c},\qquad
 \bigoplus_c\bigl(M_{d_c}^\pm\cap(V^\pm)^{\Gamma_c}\bigr),
\]
respectively. Indeed, an invariant value at a representative
extends uniquely by equivariance, independently of the element
used to transport that representative. This is the same
identification as in Proposition~\ref{acl-class-evaluation}.
Each displayed integral module is a full lattice in the
corresponding space: it is a submodule of a finite lattice, and
every vector in the invariant space has a nonzero integral
multiple in that lattice after clearing finitely many denominators.

Define
\begin{equation}\label{cp-period}
              P(f,h)=\sum_{c\in C}\frac{B(f(d_c),h(d_c))}{n_c}.
\end{equation}
Changing $d_c$ to $gd_c$ conjugates its stabilizer and preserves
its pairing by \eqref{cp-invariance}. Thus $P$ is independent
of all the representatives.

\begin{proposition}\label{cp-nondegenerate}
The pairing $P$ on $\mathcal A^-_K\times\mathcal A^+_K$
is nondegenerate. If $D$ is a positive integer divisible by
every $n_c$, then
\begin{equation}\label{cp-integrality}
       D\,P(\mathcal A^-_{\mathcal O},
                         \mathcal A^+_{\mathcal O})\subseteq\mathcal O.
\end{equation}
In fact one may take $D=\operatorname{lcm}_{c\in C}n_c$,
with the empty least common multiple defined to be $1$.
\end{proposition}
\begin{proof}
For a nonzero $x\in(V^-)^{\Gamma_c}$, choose $y\in V^+$
with $B(x,y)\ne0$. Its average
\[
                  y_0=\frac1{n_c}\sum_{\gamma\in\Gamma_c}
                                                       \rho^+(\gamma)y
\]
is invariant and satisfies $B(x,y_0)=B(x,y)$, by
\eqref{cp-invariance}. Thus the restriction to the invariant
spaces is nondegenerate on the left. The same argument on
the other side proves nondegeneracy on the right.
The direct sum of these pairings, multiplied by the nonzero
scalars $1/n_c$, remains nondegenerate. If $C$ is empty,
both spaces are zero and this statement has the same meaning.

For integral $f,h$, every $B(f(d_c),h(d_c))$ belongs to
$\mathcal O$, and $D/n_c$ is an integer. Multiplying
\eqref{cp-period} by $D$ proves \eqref{cp-integrality}.
The averaging used to prove nondegeneracy is not used to
prove this integral assertion.
\end{proof}

For a prime $w$ of $\mathcal O$, normalize $v_w$ by
$v_w(w)=1$ and put $v_w(0)=+\infty$. If $p$ is the
rational prime below $w$ and $e_w=v_w(p)$, then
\begin{equation}\label{cp-nat-valuation}
                        v_w(D)=e_w v_p(D).
\end{equation}
Indeed write $D=p^a b$ with $p\nmid b$. Bezout's identity
for $p,b$ shows that $b$ is a unit in $\mathcal O_w$.

\begin{corollary}\label{cp-local}
For integral $f,h$,
\[
                 v_w(P(f,h))\ge -e_wv_p(D).
\]
The same assertion holds if the first integral lattice is
replaced by its localization at $w$, and the second is
localized as well.
More generally, if $\mathfrak a$ is a nonzero fractional
ideal and $f\in\mathfrak a(\mathcal A^-_{\mathcal O})_w$,
$h\in(\mathcal A^+_{\mathcal O})_w$, then
\begin{equation}\label{cp-shift}
                  v_w(P(f,h))\ge v_w(\mathfrak a)-v_w(D).
\end{equation}
\end{corollary}
\begin{proof}
Localizing \eqref{cp-integrality} gives
$DP(f,h)\in\mathcal O_w$. This proves the first statements
and \eqref{cp-nat-valuation} gives their normalization.
The fractional ideal $\mathfrak a\mathcal O_w$ is principal;
write $f=a f_0$ with $v_w(a)=v_w(\mathfrak a)$ and
$f_0\in(\mathcal A^-_{\mathcal O})_w$. Bilinearity gives
\eqref{cp-shift}. Zero values cause no exception under the
chosen valuation convention.
\end{proof}

\subsection{The actual normalized polynomial pairing}

We now specify the preceding data for the row construction.
Use the dominant integral weights, determinant powers,
normalized bilinear Fock pairing $B$, and polynomial
representation $\rho_\lambda$ of Section~\ref{coef:section}.
For a tuple of full lattices $M=(M_\tau)_\tau$, put
\[
                M^+=N(M),\qquad M^-=D(M^\vee),
\]
in the notation of \eqref{coef:transport}. The integral
duality already proved in \eqref{coef:transport-duality} is
the exact equality
\begin{equation}\label{cp-polynomial-dual}
                     M^-=(M^+)^{\vee,B}.
\end{equation}
In particular $B(M^-,M^+)\subseteq\mathcal O$.
This uses the actual factorial normalization of the Fock
pairing, all determinant characters, and the unsaturated
Weyl lattice. It is valid at every coefficient prime,
including ramified primes and primes above $2$.

If the base change is $A=(A_\tau)_\tau$, the two coefficient
actions are $\rho_\lambda(A^{-t})$ and
$\rho_\lambda(A)$. Their pairing is invariant because
\[
 B(\rho_\lambda(A^{-t})x,\rho_\lambda(A)y)=B(x,y),
\]
by \eqref{coef:contravariant}. Multiplying the first action
by a finite-order level character and the second by its
inverse preserves this identity. Character values are
algebraic integral units, so they preserve the corresponding
integral lattices as well.

For the actual row objects the level-prime kernel condition
is absent: the underlying lattice is a unimodular Hermitian
lattice $N=N^\#$. The polynomial point lattice is the Weyl
lattice transported by the contragredient action on $N$;
the row lattice is the dual Weyl lattice transported
covariantly on $N$. Thus these are precisely
$D(M^\vee)$ and $N(M)$. The additional local bases in the
point description change neither classes nor stabilizers,
by Corollary~\ref{acl-row-classes}. The two level characters
are inverse, as in \eqref{acl-row-character}. Hence all
the hypotheses of Proposition~\ref{cp-nondegenerate} apply
to these actual row and point lattices.

For an embedding $\sigma:K\hookrightarrow\mathbb C$, let
$\operatorname{ev}(x_c)$ be the frame function associated
to a polynomial $x_c$ as in \eqref{acl-evaluation}.
Recovering its polynomial in the identity frame gives
exactly $\sigma(x_c)$: the determinant factor is removed
with the same exponent with which it was inserted.
The complex period in the construction is therefore
\begin{equation}\label{cp-actual-period}
 \operatorname{Per}(\operatorname{ev}(f),\operatorname{ev}(h))
      =\sum_c\frac{B_{\mathbb C}(\sigma(f(d_c)),
                                      \sigma(h(d_c)))}{n_c}
      =\sigma(P(f,h)).
\end{equation}
The last equality can be checked on monomials: the Fock
pairing is a finite sum of products of coefficients and
integer factorials, divided by the fixed nonzero rational
normalizing constant. An embedding commutes with all these
operations. No analytic comparison or unproved algebraicity
assertion is needed for this finite class sum.

Here the symbol $\operatorname{Per}$ denotes the final
polynomial period pairing. The iterated theta column is
substituted into its first argument later. Establishing
the arithmetic size of that column is a different step.

\subsection{Class restrictions and the stabilizer denominator}

For $c\in C$, let $e_c$ be restriction to the orbit of $d_c$,
followed by extension by zero. These operators preserve the
equivariant spaces, their integral lattices, and their
complexifications. They satisfy
\[
                 \sum_c e_c=1,\qquad e_ce_{c'}=0\quad(c\ne c').
\]
Give the complex point space its positive Hermitian form
\[
                   \langle x,y\rangle
                          =\sum_c n_c^{-1}\langle x_c,y_c\rangle_c,
\]
where the local forms are the normalized Fock forms. It is
conjugate-linear in its first variable. The $e_c$ are
orthogonal projections for this form, directly from the sum.

Fix a row vector $h$ and let $J(h)$ represent its period
functional, so
\[
       \langle J(h),x\rangle=\operatorname{Per}(x,\operatorname{ev}(h)).
\]
Such a vector exists uniquely because the Hermitian form is
positive definite on a finite-dimensional space. In coordinates
with Gram matrix $H$, solving this equation amounts to solving
the invertible linear system for $\overline{J(h)}^{\,t}H$.
In particular $J$ is conjugate-linear in $h$ under this
convention; no assertion of bilinearity of the Hermitian form
is being made.

\begin{proposition}\label{cp-class-contribution}
For integral $f,h$ the normalized class contribution
\[
                p_c(f,h)=n_c\langle e_cJ(h),e_c\operatorname{ev}(f)\rangle
\]
is $\sigma(B(f(d_c),h(d_c)))$, hence belongs to
$\sigma(\mathcal O)$. Moreover
\[
              \operatorname{Per}(\operatorname{ev}(f),\operatorname{ev}(h))
                           =\sum_c n_c^{-1}p_c(f,h).
\]
\end{proposition}
\begin{proof}
Orthogonality gives
$\langle e_cJ(h),e_cx\rangle=\langle J(h),e_cx\rangle$.
Apply \eqref{cp-actual-period} with $e_cf$ in place of $f$.
All terms other than $c$ vanish, leaving
$n_c^{-1}\sigma(B(f(d_c),h(d_c)))$. Multiplication by $n_c$
proves the first assertion and \eqref{cp-polynomial-dual}
proves its integrality. Summing the restrictions proves
the last identity.
\end{proof}

In the source, the row stabilizer embeds into the stabilizer
of its image in the last column class. This accounts for the
chosen common denominator. More generally, suppose a map
of objects is equivariant for an injective homomorphism of
acting groups. Its restriction to the stabilizer of an object
is injective into the stabilizer of the image, because
equivariance sends a fixing element to a fixing element.
Both groups here are finite, so Lagrange's theorem gives
divisibility of their orders. Changing the representative
of the image merely conjugates this injection.
Consequently the source's positive integer $D_\Gamma$,
divisible by the last-column stabilizer orders, is also
divisible by each $n_c$. Corollary~\ref{cp-local} gives
exactly
\begin{equation}\label{cp-source-bound}
 v_w\!\left(\sigma^{-1}\operatorname{Per}
                (\operatorname{ev}(f),\operatorname{ev}(h))\right)
                      \ge -e_wv_p(D_\Gamma).
\end{equation}
The inverse in this notation is applied only to a value
already proved to lie in $\sigma(K)$.

The same proof applies for every used finite-order twist.
If the row lattice is integral at all primes, it gives
\eqref{cp-source-bound} at all primes. If integrality has
only been supplied away from a prescribed exceptional set,
the conclusion holds away from that same set. No exceptional
prime is silently removed.

\subsection{Passage to columns and the remaining estimate}

Let $L$ be a lattice of column vectors and let
$T:L\otimes K\to\mathcal A^-_K$ be the rational map
which forms the iterated column and then its point value.
At a coefficient prime $w$, if one proves
\[
                  T(L_w)\subseteq\mathfrak a_w
                                      (\mathcal A^-_{\mathcal O})_w,
\]
then \eqref{cp-shift} immediately gives
\begin{equation}\label{cp-column-bound}
 v_w(P(Tx,h))\ge v_w(\mathfrak a_w)-v_w(D_\Gamma)
       \quad(x\in L_w,\ h\in\mathcal A^+_{\mathcal O}).
\end{equation}
This deduction retains the ramification factor in
\eqref{cp-nat-valuation}. The extension of the integral
cut and of the test-vector range is already proved in
Section~\ref{ip-section}. What remains here is the actual
containment for $T$, including the special estimates at
the level prime and the two type primes. The initial
period bound \eqref{cp-source-bound} alone does not give
that stronger containment.

For completeness, a finite set of residue-prime shifts has
an exact global interpretation. If $\mathfrak a=\prod_w w^{-b_w}$,
with integral $b_w$ of finite support, then a $K$-valued
bilinear pairing is integral on $\mathfrak a L\times M$
if and only if
\[
                     v_w(P(L,M))\ge b_w\quad\text{for every }w.
\]
Localize: $\mathfrak a_w$ has valuation $-b_w$, so the
condition is precisely that its product with every pairing
value have nonnegative valuation. An element of $K$ is
integral if and only if it is integral at every finite prime.
This proves both directions, including negative $b_w$.

Two examples show why the qualifications above are necessary.
Let a group of order two act trivially on a single object
and on $V^-=V^+=\mathbb Q$, with $B(x,y)=xy$ and both
lattices $\mathbb Z$. Then $P(1,1)=1/2$. Thus the
stabilizer denominator cannot be deleted in general, even
though the class contribution $p_c(1,1)$ is integral.
Also, over $K=\mathbb Q(\zeta_3)$, let an element act as
$\zeta_3$ on both one-dimensional coefficient spaces.
Then $B(\zeta_3x,\zeta_3y)=\zeta_3^2B(x,y)$, so a
class sum is not independent of representatives with those
two actions. Inverse characters restore invariance. Neither
example satisfies the corresponding qualified assertion
used in the submitted construction.


